% !TEX root = CoeneDylanBP.tex
\chapter{\IfLanguageName{dutch}{Stand van zaken}{State of the art}}%
\label{ch:stand-van-zaken}

% Tip: Begin elk hoofdstuk met een paragraaf inleiding die beschrijft hoe
% dit hoofdstuk past binnen het geheel van de bachelorproef. Geef in het
% bijzonder aan wat de link is met het vorige en volgende hoofdstuk.

% Pas na deze inleidende paragraaf komt de eerste sectiehoofding.

Deze literatuurstudie is opgedeeld in drie delen. In het eerste deel wordt de scoreberekening en de reglementering binnen het trampolinespringen toegelicht. Het tweede deel omvat de huidige stand van zaken op het gebied van meettechnologie{\"e}n. Het derde deel bespreekt de relevante computer vision-technieken die als bouwstenen dienen voor het systeem dat in dit onderzoek wordt ontwikkeld.

\section{Scoreberekening en Reglementering}
\label{sec:scoreberekening}

Om het belang van HD en ToF bij de beoordeling te begrijpen, wordt eerst het algemene beoordelingsproces toegelicht zoals vastgelegd door de internationale gymnastiekfederatie.

\subsection{Algemene Beoordeling}
Volgens de FIG Code of Points \parencite{FTTC2024} bestaat een individuele oefening uit een reeks van 10 sprongen, waarvan de eindscore wordt bepaald op basis van vijf componenten:

\begin{itemize}
    \item $D$ - Moeilijkheidsgraad (Difficulty)
    \item $E$ - Uitvoering (Execution)
    \item $HD$ - Horizontale verplaatsing (Horizontal Displacement)
    \item $ToF$ - Vluchttijd (Time of Flight)
    \item $P$ - Strafpunten (Penalties)
\end{itemize}

De berekening van de eindscore gebeurt aan de hand van de volgende formule:

\[
    \text{Eindscore} = E + HD + D + ToF - P
\]

In dit onderzoek ligt de focus specifiek op de HD en ToF componenten.

\subsection{Het trampolinetoestel}
\label{subsec:trampolinetoestel}

Om de HD- en ToF-metingen correct te kunnen interpreteren, is het noodzakelijk om de fysieke eigenschappen van het toestel te kennen. De FIG Code of Points \autocite{FTTC2024} legt de afmetingen en technische specificaties van het wedstrijdtoestel vast. Het volledige trampolinebed meet 5,05\,m in lengte en 3,05\,m in breedte. Het effectief bespannen gedeelte, waarop de gymnast landt, bedraagt 4,28\,m $\times$ 2,14\,m. De hoogte van het toestel bedraagt minimaal 1,15\,m boven de vloer.

Doordat deze afmetingen internationaal gestandaardiseerd zijn, is het mogelijk om in een camera-gebaseerd meetsysteem de pixelco{\"o}rdinaten van de hoekpunten van het bed te koppelen aan vaste metrische coördinaten. Dit vormt de basis voor de homografietransformatie die in dit onderzoek wordt toegepast.


\subsection{Horizontale Verplaatsing (HD)}
De HD component geeft weer hoe vaak een gymnast tijdens zijn of haar sprongreeks afweek van het midden van de trampoline. Het trampolinebed is opgedeeld in zones. Afhankelijk van de zone waarin de voeten van de atleet landen, worden er strafpunten afgetrokken van een basisscore van 10 punten \autocite{FTTC2024}.

Het midden van de trampoline geldt als ideale landingspositie zonder aftrek, terwijl elke zone die verder ligt van het midden geleidelijk hogere strafpunten oplevert. Figuur~\ref{fig:svz_HD_zones} illustreert de hoeveelheid punten die per zone worden afgetrokken van de basisscore.

\begin{figure}[ht]
    \centering
    \includegraphics[width=0.7\linewidth]{HD_zones.jpg}
    \caption{Indeling van zones op het trampolinebed voor het bepalen van strafpunten voor horizontale verplaatsing \autocite{FTTC2024}.}
    \label{fig:svz_HD_zones}
\end{figure}

\subsection{Vluchttijd (ToF)}
De ToF component omvat de totale tijd dat de gymnast zich in de lucht bevindt tijdens zijn of haar routine. Deze is belangrijk om twee redenen: enerzijds draagt ToF direct bij aan de totaalscore, anderzijds betekent een hogere ToF dat de gymnast meer tijd heeft om de technische componenten uit te voeren.

Deze tijd wordt geregistreerd in duizendsten van een seconde en vervolgens afgerond naar beneden op een honderdste van een seconde. Zo zou een gymnast met een gemeten vluchttijd van 16.197 seconden een score krijgen van 16.19 voor ToF \autocite{FTTC2024}:
\[
    16.197 \text{ seconden} \rightarrow \lfloor 16.197 \rfloor_{\downarrow 0.01} = 16.19 \text{ punten}
\]

\section{Meettechnologie{\"e}n}
\label{sec:state_of_the_art}

Er zijn verschillende methoden om de hierboven beschreven parameters te meten. Hieronder worden het huidige commerci{\"e}le meetsysteem en recent relevant onderzoek naar alternatieve camera-gebaseerde methoden besproken.

\subsection{Huidige Standaard: HDTS}
\textcite{Ferger2017} beschrijven hoe de introductie van objectieve meetwaarden voor HD en ToF heeft geleid tot de ontwikkeling van het \emph{HDTS Trampoline Measurement Device}. Dit systeem, gekoppeld aan de Eurotramp \textit{Qira 1.03} software, geldt momenteel als de standaard in competities.

Het systeem maakt gebruik van vier krachtplaten die geplaatst worden onder de poten van het trampolinerek. Voorafgaand aan gebruik kalibreert het systeem zichzelf automatisch. Op basis van de timing en verdeling van de gemeten krachten via de sensoren worden de waarden voor zowel HD als ToF berekend en real-time weergegeven \autocite{Dyas2023}.

\paragraph{Berekening Horizontale Verplaatsing}
De HD-waarde wordt bepaald door de landingsposities te berekenen op basis van de krachtverhoudingen. De coördinaten worden berekend met de volgende formules \autocite{Ferger2019}:

\begin{align*}
    F(dx) & = \frac{F_3 + F_4}{\frac{L}{2}} - \frac{F_1 + F_2}{\frac{L}{2}} \\
    F(dy) & = \frac{F_2 + F_3}{\frac{B}{2}} - \frac{F_1 + F_4}{\frac{B}{2}}
\end{align*}

Waarbij $F_{1\dots4}$ de krachten van de vier platen zijn, en $L$ en $B$ respectievelijk de lengte en breedte van de trampoline. De ori{\"e}ntatie van deze platen is weergegeven in Figuur~\ref{fig:svz_Krachtplaten}.

\begin{figure}[ht]
    \centering
    \includegraphics[width=0.7\linewidth]{Krachtplaten.jpg}
    \caption{Positie van de krachtplaten onder de trampoline en de X- en Y-richtingen voor positiebepaling \autocite{Ferger2019}.}
    \label{fig:svz_Krachtplaten}
\end{figure}

\paragraph{Berekening Vluchttijd}
ToF wordt gemeten door de krachtmetingen te combineren met een algoritme dat landings- en afstootmomenten detecteert. Het systeem zoekt naar een krachtpiek (>5000 N) en bepaalt de contactmomenten op het punt waar 50\% van die piek wordt bereikt. Vanwege de elasticiteit van het doek wordt de landing geschat op 74 ms vóór dit punt en de afstoot op 88 ms erna (zie Figuur~\ref{fig:svz_krachtcurve}) \autocite{Ferger2019}.

\begin{figure}[ht]
    \centering
    \includegraphics[width=0.5\textwidth]{ToF_meting.jpg}
    \caption{Kracht-tijd curve ter illustratie van de schatting van afstoot en landing \autocite{Ferger2019}.}
    \label{fig:svz_krachtcurve}
\end{figure}

Het HDTS-systeem biedt hoge meetnauwkeurigheid, maar kent ook praktische nadelen. Het systeem vereist een permanente installatie van vier krachtplaten onder de trampolinepoten en is daardoor enkel haalbaar in officiële wedstrijdomgevingen met de nodige infrastructuur. Voor trainingsopstellingen in clubs zonder HDTS-licentie of voor analyse buiten competitieverband is het systeem financieel en logistiek niet toegankelijk. Dit vormt de voornaamste motivatie voor de zoektocht naar een goedkoper, camera-gebaseerd alternatief.

\subsection{Computer Vision Benaderingen}
Het onderzoek van \textcite{Park2022} presenteert een systeem dat driedimensionale voetposities op een fitnesstrampoline bepaalt. Dit is relevant om de haalbaarheid van een camera-gebaseerd alternatief voor HDTS te evalueren.

Dit systeem gebruikt een camera met fisheye-lens onder de trampoline om schaduwen van de voeten te registreren. De beelden worden omgezet naar grijswaarden en genormaliseerd. Vervolgens voorspelt een deep learning-model de coördinaten.

Figuur~\ref{fig:svz_park_pipeline} toont de werking van dit algoritme: naast de visuele data wordt ook de voetgrootte als parameter meegenomen om de nauwkeurigheid van het Faster R-CNN netwerk te verhogen.

\begin{figure}[ht]
    \centering
    \includegraphics[width=0.6\linewidth]{ParkAlgoritme.jpg}
    \caption{Schematische weergave van het detectie-algoritme. De fisheye-beelden worden omgezet naar grijswaarden en samen met de voetgrootte verwerkt door een Faster R-CNN netwerk om de XYZ-coördinaten te bepalen \autocite{Park2022}.}
    \label{fig:svz_park_pipeline}
\end{figure}

De resultaten van dit onderzoek tonen een gemiddelde afwijking van 11.7 mm. Dit suggereert dat deze methode, mits aangepast voor grote trampolines, voldoende nauwkeurig kan zijn voor het bepalen van de landingszones. De methode is echter specifiek ontwikkeld voor fitnesstrampolines en vereist een cameraopstelling onder het toestel, wat in een wedstrijdcontext moeilijk realiseerbaar is wegens toegankelijkheidsbeperkingen en veiligheidsvoorschriften. Geïnspireerd door deze resultaten werd in de initi{\"e}le fase van dit onderzoek eveneens ge{\"e}xperimenteerd met een cameraopstelling onder de trampoline. Dit wordt verder toegelicht in Hoofdstuk~\ref{ch:methodologie}.

\section{Computer vision-technieken}
\label{sec:computer-vision-technieken}

De volgende subsecties bespreken de computer vision-technieken die als bouwstenen dienen voor de ontwikkelde pipeline. Hierbij komen positiebepaling via homografie, objectdetectie, pose estimation en beeldsegmentatie aan bod, gevolgd door een bespreking van de trainingsdata.

% Beschrijving van hoe hawk-eye een inspiratie was voor homografie-gebaseerde positiebepaling.
\subsection{Positiebepaling met Homografie}
\label{subsec:hawk_eye}

Naast de methode die gebruikmaakt van het registreren van schaduwen, is het ook mogelijk om camerabeelden vanaf een vast standpunt in de ruimte te kalibreren ten opzichte van een gegeven veld. Het bepalen van een positie van een persoon of object met een camera door middel van een homografie-transformatie is sterk gebaseerd op het bekende Hawk-Eye-systeem, dat wordt gebruikt in zowel professionele tennis- als voetbalwedstrijden.

Recent onderzoek, zoals beschreven door \textcite{Manikanta2025}, illustreert de werking van een geautomatiseerd systeem voor tennisanalyse. In hun werk worden deep learning-modellen (zoals YOLOv5, YOLOv8 en ResNet50) geïntegreerd om spelers, de bal en de keypoints van het veld te detecteren. Door middel van de gedetecteerde keypoints kan een homografie-matrix worden opgesteld. Deze matrix maakt het mogelijk om de cameraprojectie te transformeren naar een gestandaardiseerd bovenaanzicht, wat precieze positiebepaling van acties of atleten mogelijk maakt, onafhankelijk van het perspectief van de camera. Figuur~\ref{fig:tennis-keypoints-minicourt} illustreert dit principe. De gedetecteerde keypoints van het veld worden gebruikt om de spelersposities te projecteren op een vereenvoudigde 2D-kaart.

\begin{figure}[h]
    \centering
    \includegraphics[width=0.8\textwidth]{tennis_keypoints_minicourt}
    \caption{Detectie van veldkeypoints en homografie-transformatie naar een 2D-minikaart in een tennisanalysesysteem \autocite{Kosolapov2024}.}
    \label{fig:tennis-keypoints-minicourt}
\end{figure}

Concreet betekent dit dat elk punt in het camerabeeld, uitgedrukt in pixels, eenduidig kan worden omgezet naar een punt in een tweedimensionaal vlak met correcte verhoudingen, op voorwaarde dat de vier referentiepunten nauwkeurig bekend zijn.

In dit onderzoek wordt een gelijkaardige methodiek gehanteerd om de prestaties bij het trampolinespringen te analyseren met behulp van slechts één camera. Door de homografie-transformatie toe te passen op het trampolinebed, zou de positie van de gymnast bij elke landing kunnen worden geprojecteerd. Op deze manier trachten we de horizontale verplaatsing op een objectieve en geautomatiseerde wijze te meten, waarbij de optische vervorming van de camera wiskundig wordt weggewerkt om de exacte coördinaten op het trampolinebed te bepalen.

\subsection{Objectdetectie: YOLO}
\label{subsec:yolo}

Een centrale uitdaging in het voorgestelde systeem is het per frame lokaliseren van zowel de trampoline als de gymnast. Hiervoor wordt gebruikgemaakt van de YOLO-familie (\textit{You Only Look Once}) van objectdetectiemodellen \autocite{Jocher_Ultralytics_YOLO_2023}.

In tegenstelling tot twee-staps detectoren, zoals Faster R-CNN dat door \textcite{Park2022} werd ingezet, voert YOLO de volledige detectie in één enkele doorgang door het neurale netwerk uit. Het beeld wordt opgedeeld in een raster, waarbij elke cel tegelijk een bounding box, een betrouwbaarheidsscore en een klasse-indeling voorspelt. Dit maakt het model bijzonder geschikt voor real-time toepassingen. De meest recente generatie, YOLOv11 van Ultralytics, ondersteunt naast standaard objectdetectie ook een pose-variant voor keypointdetectie en een \textit{classificatie}-variant.

% TODO: Figuur van yolo proces toevoegen
% TODO: Figuur object detectie + pose detectie?

In dit onderzoek wordt YOLOv11 ingezet voor drie afzonderlijke taken: het detecteren van de vier hoekpunten van het trampolinebed (via de pose-variant), het classificeren van frames als \textit{Landed} of \textit{Not Landed}, en het lokaliseren van de gymnast als persoon. De eerste twee modellen worden daarvoor specifiek gefinetuned op trampolinedata, zoals besproken in Sectie~\ref{sec:datasets}. De veelzijdigheid van YOLOv11, die detectie, classificatie en keypointdetectie combineert in één modelfamilie, maakt het bouwen van een geïntegreerde pipeline zeer efficiënt, waarbij meerdere deeltaken met hetzelfde framework worden opgelost.

\subsection{Pose Estimation}
\label{subsec:pose-estimation}

Pose Estimation is een computervisietechniek waarbij de posities van specifieke lichaamspunten (\textit{keypoints}) van een persoon worden bepaald op basis van een camerabeeld. In het meestgebruikte formaat; het COCO-formaat (Common Objects in Context), worden 17 keypoints gedefinieerd en gedetecteerd, waaronder de enkels, knieën, heupen, schouders en het hoofd. Deze keypoints kunnen worden gebruikt om de houding en beweging van een persoon te analyseren.

Codefragment~\ref{lst:coco-keypoints} illustreert hoe een COCO-annotatie er in de praktijk uitziet. Elk keypoint wordt opgeslagen als een triplet \texttt{(x,~y,~v)}: de pixelco{\"o}rdinaten gevolgd door een zichtbaarheidswaarde \texttt{v}. De waarde \texttt{2} geeft aan dat het keypoint zichtbaar is in het beeld; \texttt{1} betekent dat het aanwezig is maar wordt afgeschermd (\textit{occluded}); \texttt{0} betekent dat het keypoint niet gelabeld werd.

\begin{listing}[ht]
    \begin{minted}{javascript}
{
  "keypoints": [
    "nose",          "left_eye",      "right_eye",
    "left_ear",      "right_ear",     "left_shoulder",
    "right_shoulder","left_elbow",    "right_elbow",
    "left_wrist",    "right_wrist",   "left_hip",
    "right_hip",     "left_knee",     "right_knee",
    "left_ankle",    "right_ankle"
  ],
  "annotations": [{
    "keypoints": [
      320, 112, 2,  /* nose         (x, y, v=2: zichtbaar)  */
      308, 103, 2,  /* left_eye                              */
      332, 103, 2,  /* right_eye                             */
      /* ... 11 andere keypoints ...                         */
      278, 384, 2,  /* left_ankle                            */
      341, 381, 2   /* right_ankle                           */
    ],
    "num_keypoints": 17
  }]
}
\end{minted}
    \caption[COCO keypointannotatie]{Structuur van een COCO-annotatie. Elk keypoint bestaat uit drie waarden \texttt{(x,~y,~v)}: pixelco{\"o}rdinaten en een zichtbaarheidswaarde \texttt{v} (0~=~niet gelabeld, 1~=~occluded, 2~=~zichtbaar). De volledige array telt 51 waarden (17 keypoints $\times$ 3).}
    \label{lst:coco-keypoints}
\end{listing}

ViTPose \autocite{Xu2022} is een state-of-the-art model voor keypointdetectie dat gebruik maakt van een \textit{Vision Transformer} (ViT) als onderliggend basisnetwerk (backbone). Transformerarchitecturen zijn oorspronkelijk ontwikkeld voor tekstverwerking, maar zijn later succesvol toegepast op beeldtaken doordat ze globale verbanden in een beeld kunnen modelleren via een aandachtsmechanisme (attention mechanism), waarbij het model zelf bepaalt welke beeldregio's het meest relevant zijn voor de taak. ViTPose behaalt topprestaties op de COCO keypointbenchmarks en is daarmee een geschikte kandidaat voor het nauwkeurig lokaliseren van lichaamspunten.

% TODO: Figuur van ViT Attention mechanisme toevoegen

De geschiktheid van ViTPose voor trampolinecontext wordt bevestigd door recent onderzoek van \textcite{DroletRoy2026}, die het model gebruikten om verder te finetunen op synthetische trampolinedata en daarmee \textit{state-of-the-art} 2D-keypointprestaties behalen op echte trampolineopnames. Hun resultaten tonen aan dat een gefinetuned ViTPose-model ook bij 3D-reconstructie van beelden uit meerdere camerastandpunten beter presteert dan het voorgetrainde basismodel. Er werd een gemiddelde reductie van 19,6\% in Mean Per Joint Position Error (MPJPE) gerapporteerd. Dit benadrukt het potentieel van ViTPose voor de uitdagende poses en gezichtshoeken die in trampolinegymnastiek voorkomen.

% TODO: Figuur van paper van DroletRoy toevoegen???

Voor dit onderzoek was de initiële aanpak om het middelpunt tussen de twee enkels (keypoints 16 en 17) te gebruiken als nauwkeurig referentiepunt voor de landingspositie. In de praktijk bleek echter dat bij snelle trampolinesprongen de bewegingsonscherpte (\textit{motion blur}) in het landingsframe regelmatig leidt tot het wegvallen van keypointpredities. Omdat dit zich specifiek voordoet op het kritieke moment van de landing, is deze aanpak onbetrouwbaar voor dit toepassingsdomein.

\subsection{Segmentatie: SAM-2}
\label{subsec:sam2}

Als alternatief voor keypointdetectie wordt in dit onderzoek gebruikgemaakt van pixel-niveau segmentatie. In plaats van een bounding box rondom de gymnast te berekenen, produceert een segmentatiemodel een binair masker dat exact aangeeft welke pixels tot het object behoren.

Segment Anything Model 2 (\textit{SAM-2}) \autocite{Ravi2024} is een algemeen segmentatiemodel ontwikkeld door Meta, ontworpen voor zowel statische beelden als videosequenties. Wat SAM-2 kenmerkt is de prompt-gestuurde aanpak. Het model accepteert een eenvoudige invoer, zoals een bounding box of een pixelpunt, en genereert op basis daarvan een nauwkeurig segmentatiemasker zonder specifieke finetuning. Figuur~\ref{fig:sam2-prompt} illustreert hoe een combinatie van een bounding box en positieve of negatieve punten als prompt kan worden meegegeven om het gewenste object te segmenteren.

\begin{figure}[ht]
    \centering
    \includegraphics[width=0.85\textwidth]{sam_prompt_based.png}
    \caption[Prompt-gestuurde segmentatie met SAM-2]{Voorbeeld van prompt-gestuurde segmentatie met SAM-2. Een bounding box en optionele positieve of negatieve punten worden als invoer meegegeven. SAM-2 genereert op basis hiervan een pixelnauwkeurig segmentatiemasker van het gewenste object \autocite{Ravi2024}.}
    \label{fig:sam2-prompt}
\end{figure}

Dit biedt een belangrijk voordeel ten opzichte van pose estimation in de context van trampolinespringen. Doordat de laagste pixel van het segmentatiemasker direct wordt geëxtraheerd als voetcontactpunt, is er geen afhankelijkheid van expliciete keypointvoorspellingen. Dit maakt de aanpak robuuster voor situaties met veel motion blur. In de voorgestelde pipeline levert de persoondetectie via YOLOv11n een bounding box die als prompt aan SAM-2-tiny wordt meegegeven. Het laagste punt van het resulterende masker dient vervolgens als landingspixel voor de homografietransformatie. De gekozen variant SAM-2-tiny biedt een gunstige balans tussen segmentatienauwkeurigheid en berekentijd, wat relevant is voor toepassingen waarbij meerdere frames snel na elkaar moeten worden verwerkt.

\subsection{Optical Flow en de Farneback-methode}
\label{subsec:farneback}

Optical flow beschrijft de schijnbare beweging van pixels tussen twee opeenvolgende videoframes. Het basisidee is intuïtief: als een object verschuift, verschuift ook het bijbehorende patroon van lichtintensiteiten in het beeld. Figuur~\ref{fig:optical-flow-concept} illustreert dit principe: een beeldkenmerk op positie $(x, y)$ op tijdstip $t$ is terug te vinden op positie $(x + \mathrm{d}x,\, y + \mathrm{d}y)$ op tijdstip $t + \mathrm{d}t$.

\begin{figure}[ht]
    \centering
    \includegraphics[width=0.6\linewidth]{optical_flow_intuitief.png}
    \caption[Basisprincipe van optical flow]{Het basisprincipe van optical flow: een beeldkenmerk op positie $(x,y)$ op tijdstip $t$ wordt teruggevonden op positie $(x+\mathrm{d}x,\, y+\mathrm{d}y)$ op tijdstip $t+\mathrm{d}t$.}
    \label{fig:optical-flow-concept}
\end{figure}

Traditionele methoden zoals Lucas-Kanade gaan ervan uit dat de beweging constant is binnen een kleine omgeving van pixels: alle naburige pixels bewegen gelijk. Wiskundig is dit vergelijkbaar met het benaderen van de lokale beeldstructuur als een plat vlak. Dit werkt voor kleine, trage verplaatsingen, maar schiet tekort bij complexe of rijke beeldstructuren.

\textcite{Farneback2003} introduceert een fundamenteel andere benadering. In plaats van de lokale omgeving van elk pixel als een vlak te beschouwen, modelleert hij die omgeving als een gekromd oppervlak. Stel je voor dat de helderheidswaarden rondom een pixel een driedimensionaal berglandschap vormen: klassieke methoden leggen daar een plat vlak overheen, terwijl Farneback een gebogen vorm toepast die de lokale structuur veel nauwkeuriger beschrijft. Door te vergelijken hoe dit gekromde oppervlak verschuift tussen twee opeenvolgende frames, berekent het algoritme voor elke pixel afzonderlijk een verplaatsingsvector. Het resultaat is een \textit{dense flow field} (hoge dichtheid): een volledig beeld van beweging, pixel voor pixel, in tegenstelling tot sparse (lage dichtheid) methoden die enkel voor een selectie van kenmerkpunten werken. Dit maakt Farneback geschikt om subtiele, gedistribueerde vervormingen in een beeld te detecteren. De methode is beschikbaar als \texttt{cv2.calcOpticalFlowFarneback()} in de OpenCV-bibliotheek \autocite{opencv_library}.

\subsection{Datasets}
\label{sec:datasets}

Het trainen van domeinspecifieke computervisiemodellen vereist data die representatief is voor de omgeving waarin ze zullen worden gebruikt. Voor eerder niche sporten zoals trampolinespringen zijn dergelijke datasets schaars, wat de ontwikkeling van geautomatiseerde meetsystemen een stuk moeilijker maakt.

\subsubsection{Trampolinedataset}
\label{subsec:dataset-lin}

\textcite{Lin2025} stellen een dataset voor, specifiek ontworpen voor de beoordeling van actiekwaliteit in trampolinegymnastiek. De dataset bevat 206 video-opnames met finale HD- en ToF-score annotaties van competitieve trampolinesprongen. De dataset is gepubliceerd in het kader van onderzoek naar een causaal cross-modal framework voor langetermijn beoordelingstaken. De dataset bevat data van vier competities: Buenos Aires 2018, Tokyo 2020, Hangzhou 2022 en Paris 2024.

% TODO: Figuur voorbeeldframes uit de dataset toevoegen

In dit onderzoek wordt de dataset opgesplitst, waarbij een deel wordt gebruikt als trainingsdata voor twee modellen uit de pipeline: het finetunen van het hoekpuntdetectiemodel dat de vier hoeken van het trampolinebed lokaliseert, en het trainen van de landingsclassifier die per frame bepaalt of de gymnast al dan niet in contact is met het doek. Daarnaast wordt een apart gehouden deel van de dataset ingezet als testset voor het evalueren van de nauwkeurigheid van het volledige systeem, door de gemeten HD- en ToF-waarden te vergelijken met de werkelijke scores. De variatie in camerastandpunten en competitieomgevingen binnen de dataset maakt het een geschikte bron voor het trainen van modellen die robuust moeten zijn over verschillende opstellingen.